\chapter{Нейрони грають у теніс}
% full version: chapters/22_dishbrain.tex

Найвідоміший дослід із живими нейронами на чипі --- DishBrain, «мозок у чашці». Близько мільйона нейронів, вирощених на чипі з двадцятьма шістьма тисячами електродів, грали в найпростіший комп'ютерний теніс: м'яч літає екраном, ракетка ходить угору й униз.

\section*{Як влаштована гра}

М'яч потрапляє в мережу через вісім електродів. Який із них клацає --- показує, на якій висоті м'яч. Як часто клацає --- показує, чи далеко м'яч: чотири рази на секунду, коли він біля дальньої стінки, і дедалі частіше, до сорока, коли він підлітає до ракетки. Вихід --- дві ділянки електродів: клацання в одній рухають ракетку вгору, в іншій --- униз. Кожні десять мілісекунд комп'ютер рахує клацання й рухає ракетку. Ці десять мілісекунд --- такт комп'ютера стенда, а не культури: сама мережа працює асинхронно, як будь-яка нервова мережа.

\pencilfig{22}{\pencilart{22}}{М'яч, ракетка й мільйон нейронів, які її рухають.}

\section*{Учитель}

Жодного дофаміну. Після вдалого удару мережа отримує короткий і цілком передбачуваний сигнал: усі вісім електродів разом, одна десята секунди. Після промаху --- чотири секунди безладної стимуляції у випадкових місцях і у випадкові моменти, потім пауза, і нова подача.

\section*{Що вийшло}

За двадцятихвилинний сеанс розіграші в живих культур ставали довшими, а пропущених подач ставало менше. Найкраще це виходило за повного зворотного зв'язку. Коли замість сигналів після удару й промаху просто наставала тиша, поліпшення теж було, але слабше; без зворотного зв'язку його не було. Людські нейрони грали трохи краще за мишачі. А от навчання, розтягнуте на кілька днів, надійним не виявилося.

\section*{Чому мережа вчиться}

Автори пояснюють результат так: мережа прагне зробити свої відчуття передбачуваними. Промах приносить хаос, удар --- порядок, і мережа змінюється, щоб порядку стало більше. Але є й приземленіше пояснення, яке теж узгоджується з даними: «струшуй після промаху, не чіпай після удару». Кожен струс трохи перемішує зв'язки навмання. Налаштування, за яких мережа часто хибить, весь час перемішуються, а вдалі лишаються незайманими. Мережа сама собою затримується там, де промахів мало, --- ніхто її не вчив, її просто перестають трусити. У нашій простій моделі таке правило справді піднімає частку ударів.

\pencilfig{130}{%
% scrambler: miss probability p over one setting of the network (valley in the middle); a miss kicks the setting
% at random, a hit leaves it alone, so the time spent at a setting is proportional to 1/p (6x more in the valley here)
\draw[pen] (0,1.2) -- (8.2,1.2);
\draw[penlite=pcRed] plot[smooth] coordinates {(0.0,2.46) (0.8,2.46) (1.6,2.46) (2.0,2.44) (2.4,2.41) (2.8,2.32) (3.2,2.15) (3.6,1.90) (4.0,1.62) (4.4,1.44) (4.8,1.44) (5.2,1.62) (5.6,1.90) (6.0,2.15) (6.4,2.32) (6.8,2.41) (7.2,2.44) (7.6,2.46) (8.0,2.46)};
\node[lbl, text=pcRed, anchor=south] at (7.55,2.55) {промахів багато};
\node[lbl, text=pcRed, anchor=west] at (5.3,1.45) {мало};
% kicked balls on the slopes
\draw[dotpen=pcBlue, wash=pcBlue] (1.3,2.68) circle (0.12);
\draw[arr=pcGray, densely dashed] (1.35,2.82) to[bend left=50] (2.35,2.72);
\draw[arr=pcGray, densely dashed] (1.22,2.82) to[bend right=50] (0.4,2.72);
\draw[dotpen=pcBlue, wash=pcBlue] (6.3,2.45) circle (0.12);
\draw[arr=pcGray, densely dashed] (6.25,2.59) to[bend right=50] (5.45,2.25);
\node[lbl, anchor=south] at (1.3,3.05) {струс після промаху};
% the ball left alone in the valley
\draw[dotpen=pcBlue, wash=pcBlue] (4.6,1.66) circle (0.12);
\node[lbl, anchor=south] at (4.6,1.95) {не трусять};
% where the network spends its time (proportional to 1/p)
\fill[pcGreen, opacity=0.22] (0,0) -- plot[smooth] coordinates {(0.0,0.18) (0.8,0.18) (1.6,0.18) (2.0,0.19) (2.4,0.19) (2.8,0.21) (3.2,0.24) (3.6,0.33) (4.0,0.55) (4.4,0.98) (4.8,0.98) (5.2,0.55) (5.6,0.33) (6.0,0.24) (6.4,0.21) (6.8,0.19) (7.2,0.19) (7.6,0.18) (8.0,0.18)} -- (8.0,0) -- cycle;
\draw[penlite=pcGreen] plot[smooth] coordinates {(0.0,0.18) (0.8,0.18) (1.6,0.18) (2.0,0.19) (2.4,0.19) (2.8,0.21) (3.2,0.24) (3.6,0.33) (4.0,0.55) (4.4,0.98) (4.8,0.98) (5.2,0.55) (5.6,0.33) (6.0,0.24) (6.4,0.21) (6.8,0.19) (7.2,0.19) (7.6,0.18) (8.0,0.18)};
\draw[pen] (0,0) -- (8.2,0);
\node[lbl, text=pcGreen!60!black, anchor=west] at (5.3,0.6) {де мережа проводить час};
\node[lbl, anchor=north] at (4.1,-0.03) {налаштування мережі};
}{Промах струшує налаштування мережі навмання, удар лишає їх у спокої. Там, де промахів багато, мережа не затримується. Де їх мало, її перестають трусити, і вона лишається там --- хоча ніхто її туди не вів.}

Як розрізнити два пояснення? Потрібен дослід, якого ще не було: після промаху подавати не хаотичний, а \emph{передбачуваний} сигнал тієї самої тривалості й сили. Якщо мережа вчиться й так, річ не в передбачуваності, а просто в різниці між «після удару» і «після промаху». Цей дослід ще чекає на свого виконавця.

Ось і кінець нашого шляху. Машина на двадцяти ватах зібрана з простих деталей, але як саме вона складає їх у думку, ми поки розуміємо лише частково. Можливо, наступний розділ напишете ви.

\fullversion{22}
